\documentclass[11pt]{article}
\usepackage[margin=1in]{geometry}
\usepackage{xcolor}
\usepackage{longtable}
\usepackage{enumitem}
\usepackage{parskip}
\usepackage{hyperref}
\usepackage{titlesec}
\definecolor{passgreen}{HTML}{1B7A3D}
\definecolor{partialorange}{HTML}{B5701A}
\definecolor{failred}{HTML}{B22222}
\definecolor{panelgray}{HTML}{F0F0EE}
\definecolor{navy}{HTML}{1F2A44}
\hypersetup{colorlinks=true,linkcolor=navy,urlcolor=navy}
\titleformat{\section}{\large\bfseries\color{navy}}{}{0pt}{}
\setlength{\parskip}{4pt}
\newcommand{\fieldlabel}[1]{\textbf{\color{navy}#1}\quad}
\begin{document}

\begin{center}
{\LARGE\bfseries\color{navy} polilabs Agent Evaluation}\\[6pt]
{\large 20-Item Quality \& Behavior Suite}\\[10pt]
\normalsize Model: claude-sonnet-4-6 \quad|\quad Run: 2026-05-22T11:31:32 \quad|\quad Suite wall time: 516.7\,s
\end{center}
\vspace{6pt}\hrule\vspace{10pt}
\section*{Summary}
\noindent\fcolorbox{navy}{panelgray}{\parbox{0.97\textwidth}{\textbf{Score:} \textcolor{passgreen}{\textbf{14 PASS}} \;/\; \textcolor{partialorange}{\textbf{4 PARTIAL}} \;/\; \textcolor{failred}{\textbf{2 FAIL}} \quad out of 20 items \quad(70\% full pass).}}

\subsection*{Key findings}
The agent scored 14 PASS / 4 PARTIAL / 2 FAIL across the 20-item suite (70\% full pass). The result is encouraging where it matters most and clear about where it does not.

Strength --- anti-hallucination holds. Every probe designed to bait fabrication was handled cleanly. Out-of-corpus state law (Q15), foreign law (Q16), a nonexistent bill (Q17), and an outright false premise (Q20, "the AI Safety Act of 2022") all produced honest refusals that named the scope boundary and declined to answer from training data. The honest-null probe (Q10, "frontier model") correctly returned "no bills define this" without inventing a definition. This is the core product thesis and it is working.

Strength --- substantive retrieval is solid. Single-bill summaries (Q3, Q11), the definitions aggregate (Q8), the amendment-graph question (Q9), and the long comparison (Q12) were accurate, citation-grounded, and faithful to the corpus, including honest caveats about extraction gaps (Q9 flagged bills whose amendment operations were not parsed). Multi-turn memory mechanics work: Q13 turn 2 resolved "it" to the prior bill and answered from context with zero new tool calls.

Weakness 1 --- no aggregate / enumeration primitive (Q6, Q7, Q18, and partly Q19). The agent cannot count or group the corpus by Congress or sponsor. It behaves honestly under this gap --- it hedges rather than fabricates --- but it cannot deliver clean answers to "how many" or "who sponsored the most." Q18 is the sharp symptom: the same question ("how many 118th-Congress bills?") asked four times produced four materially different answers --- one confidently correct (100), three evasive, two of which wrongly speculated the true count is higher than 100. The question has one true answer; the agent gave four. This is the dominant finding and the fix is a corpus count/filter/group-by tool.

Weakness 2 --- keyword-retrieval gap (Q5). Asked which bills restrict advanced-chip exports, the agent concluded the corpus contains "no standalone bills" on the topic --- when it in fact contains at least three (the CLOUD AI Act, the GAIN AI Act, and S. 321). BM25 keyword search did not surface them under the query phrasing and the agent gave up. It did not fabricate, but "honest but wrong" is still wrong. This is the concrete case for hybrid semantic retrieval.

Behavior flag --- Q14 turn 1. The agent characterized three bills' facial-recognition provisions without pulling their text, then in turn 2 opened with "I need to correct myself --- the previous response was partially fabricated." The self-correction is a good sign, but the agent should not describe a bill's provisions before reading it. Tighten the system prompt to forbid characterizing un-retrieved bills.

Bottom line. The agent fails safely, not dangerously: it hedges and returns honest nulls rather than hallucinating. For a legislative-research tool that is the right failure mode. But Q5 and Q18 show that safe failure is still failure for the user. The two priorities are unambiguous: (1) add a corpus aggregate/count tool; (2) add hybrid semantic retrieval. Neither requires touching the anti-hallucination design, which is sound.

Scope note: this run executed against the pre-streaming master agent; the streaming change affects delivery latency only, not answer content, so these quality results stand for both.
\subsection*{Scorecard}
\begin{longtable}{p{1.3cm} p{8.4cm} p{2.6cm}}
\textbf{Item} & \textbf{Category} & \textbf{Grade}\\ \hline
\endhead
Q1 & corpus scope & \colorbox{passgreen}{\textcolor{white}{\bfseries PASS}}\\
Q2 & single-bill fact & \colorbox{passgreen}{\textcolor{white}{\bfseries PASS}}\\
Q3 & single-bill substantive & \colorbox{passgreen}{\textcolor{white}{\bfseries PASS}}\\
Q4 & conceptual retrieval & \colorbox{passgreen}{\textcolor{white}{\bfseries PASS}}\\
Q5 & conceptual retrieval (vocabulary gap) & \colorbox{failred}{\textcolor{white}{\bfseries FAIL}}\\
Q6 & aggregate count & \colorbox{partialorange}{\textcolor{white}{\bfseries PARTIAL}}\\
Q7 & aggregate fact & \colorbox{partialorange}{\textcolor{white}{\bfseries PARTIAL}}\\
Q8 & definitions aggregate & \colorbox{passgreen}{\textcolor{white}{\bfseries PASS}}\\
Q9 & amendment graph & \colorbox{passgreen}{\textcolor{white}{\bfseries PASS}}\\
Q10 & definitions (honest-null probe) & \colorbox{passgreen}{\textcolor{white}{\bfseries PASS}}\\
Q11 & single-bill substantive & \colorbox{passgreen}{\textcolor{white}{\bfseries PASS}}\\
Q12 & long comparison & \colorbox{passgreen}{\textcolor{white}{\bfseries PASS}}\\
Q13 & session memory (multi-turn) & \colorbox{passgreen}{\textcolor{white}{\bfseries PASS}}\\
Q14 & session memory (multi-turn) & \colorbox{partialorange}{\textcolor{white}{\bfseries PARTIAL}}\\
Q15 & out-of-corpus (state law) & \colorbox{passgreen}{\textcolor{white}{\bfseries PASS}}\\
Q16 & out-of-corpus (foreign law) & \colorbox{passgreen}{\textcolor{white}{\bfseries PASS}}\\
Q17 & out-of-corpus (nonexistent bill) & \colorbox{passgreen}{\textcolor{white}{\bfseries PASS}}\\
Q18 & consistency (repeated question) & \colorbox{failred}{\textcolor{white}{\bfseries FAIL}}\\
Q19 & long analytical & \colorbox{partialorange}{\textcolor{white}{\bfseries PARTIAL}}\\
Q20 & false premise (hallucination probe) & \colorbox{passgreen}{\textcolor{white}{\bfseries PASS}}\\
\end{longtable}
\newpage
\section*{Item detail}
\subsection*{Q1 \quad \colorbox{passgreen}{\textcolor{white}{\bfseries PASS}}}
\fieldlabel{Category.} corpus scope\par
\fieldlabel{Question.} How many bills are in this corpus, and which Congresses and time span do they cover?\par
\fieldlabel{Expected (corpus ground truth).}\par
191 bills, covering the 118th Congress (100 bills) and 119th Congress (91 bills); introduction dates span January 9, 2023 to May 4, 2026.

\fieldlabel{Agent answer.}\quad{\small\itshape 5.5s, ttft 5.5s, tools: corpus\_coverage}\par
\textbf{191 bills} total, spanning \textbf{January 9, 2023 -- May 4, 2026}, across the \textbf{118th and 119th Congresses}.

Breakdown by tier: \\
- \textbf{Tier A} (primary/substantial AI focus): 184 bills \\
- \textbf{Tier B} (partial/peripheral AI relevance): 7 bills

Sources are drawn from GovInfo and Congress.gov, last fetched May 17, 2026. The corpus covers \textbf{legislation only} --- regulatory actions and executive orders are explicitly out of scope for v1.

\fieldlabel{Assessment.} \colorbox{passgreen}{\textcolor{white}{\bfseries PASS}}\par
Correct on every checkable point: 191 bills, 118th + 119th Congress, span January 9 2023 to May 4 2026, plus the 184/7 tier split. It did not give the 100/91 per-Congress split, but that is a tool-coverage gap (corpus\_coverage does not expose it) rather than an agent error, and the split was not strictly asked for.
\vspace{4pt}\hrule\vspace{8pt}
\subsection*{Q2 \quad \colorbox{passgreen}{\textcolor{white}{\bfseries PASS}}}
\fieldlabel{Category.} single-bill fact\par
\fieldlabel{Question.} Who sponsored the AI Fraud Deterrence Act, and when was it introduced?\par
\fieldlabel{Expected (corpus ground truth).}\par
The AI Fraud Deterrence Act (H.R. 10125, 118th Congress) was sponsored by Rep. Ted Lieu (D-CA-36) and introduced on November 14, 2024.

\fieldlabel{Agent answer.}\quad{\small\itshape 5.22s, ttft 5.21s, tools: search\_corpus}\par
The \textbf{AI Fraud Deterrence Act} has been introduced twice by the same sponsor:

- \textbf{118th Congress} --- H.R. 10125, sponsored by \textbf{Rep. Ted Lieu (D-CA-36)}, introduced \textbf{November 14, 2024}. \\
- \textbf{119th Congress} --- H.R. 6306, sponsored by \textbf{Rep. Ted Lieu (D-CA-36)}, introduced \textbf{November 25, 2025}.

Rep. Lieu reintroduced the bill in the new Congress with essentially the same title.

\fieldlabel{Assessment.} \colorbox{passgreen}{\textcolor{white}{\bfseries PASS}}\par
Sponsor Rep. Ted Lieu (D-CA-36) and introduction date November 14, 2024 both correct. It also surfaced a 119th-Congress reintroduction (H.R. 6306) --- extra, consistent detail beyond the expected answer.
\vspace{4pt}\hrule\vspace{8pt}
\subsection*{Q3 \quad \colorbox{passgreen}{\textcolor{white}{\bfseries PASS}}}
\fieldlabel{Category.} single-bill substantive\par
\fieldlabel{Question.} What does the AI Labeling Act of 2023 require, and who must comply with it?\par
\fieldlabel{Expected (corpus ground truth).}\par
Exists as House (H.R. 6466) and Senate (S. 2691) companion bills, identical text. Requires generative AI systems to attach clear, conspicuous, permanent disclosures that content is AI-generated (metadata for image/video/audio, notice for text/chatbots). Developers and third-party licensees must prevent downstream removal of disclosures. FTC enforces as an unfair/deceptive act.

\fieldlabel{Agent answer.}\quad{\small\itshape 26.32s, ttft 4.73s, tools: search\_corpus, get\_bill, get\_bill, get\_section, get\_section}\par
There are two versions --- a House bill (H.R. 6466) and a Senate companion (S. 2691). Let me pull both simultaneously.Both bills have an identical structure. Let me pull the substantive sections from both simultaneously.The two bills --- S. 2691 (Schatz, introduced July 2023) and H.R. 6466 (Kean, introduced November 2023) --- are textually identical. Here is a full breakdown drawn from Sec. 2 of S. 2691, 118th Cong. and Sec. 2 of H.R. 6466, 118th Cong.

\textbf{Who Must Comply}

Three categories of actors bear distinct duties:

1. \textbf{Generative AI systems themselves} (i.e., the operator of any system using interstate/foreign commerce to produce content) --- the primary labeling obligation sits at the system level. \\
2. \textbf{Developers of generative AI systems} --- must implement reasonable procedures to prevent downstream removal of disclosures. \\
3. \textbf{Third-party licensees} of generative AI systems --- same downstream-prevention obligation as developers, applied to their own users.

\textbf{Core Requirements}

\textbf{Labeling AI-Generated Content}

| Content Type | What the system must do | \\
|---|---| \\
| \textbf{Image, video, audio, or multimedia} | Attach a *clear and conspicuous* notice identifying the content as AI-generated; embed in metadata the identity of the tool used and the date/time of creation; make the disclosure permanent/non-removable to the extent technically feasible. | \\
| \textbf{Text} (including chatbot output) | Include a *clear and conspicuous* disclosure identifying the content as AI-generated; make it permanent/non-removable to the extent technically feasible. |

"\textbf{Clear and conspicuous}" is precisely defined: for visual content it must stand out by size, contrast, and location; for audio it must be at a volume and cadence a reasonable person can easily hear; it must not be avoidable, must use plain language, and must not be contradicted elsewhere in the communication.

"\textbf{AI-generated content}" is content *substantially* created or modified by a generative AI system such that use of the system *materially alters* the meaning or significance a reasonable person would take from it --- incidental edits are not covered.

\textbf{Developer \& Licensee Obligations}

Both developers and third-party licensees must: \\
- \textbf{Contract} with downstream users/licensees requiring them not to remove disclosures. \\
- \textbf{Require certification} that users will comply. \\
- \textbf{Terminate access} when they have reason to believe a user has removed a required disclosure.

\textbf{Government Role}

\textbf{NIST} must, within 90 days of enactment, convene an *AI-Generated Content Consumer Transparency Working Group* --- including federal agencies, AI developers, content-detection standards bodies, social media platforms, and academics. Within 1 year, the working group must publish a framework of: \\
- Technical standards for AI-content detection (image, video, audio, multimedia). \\
- Best practices for platforms to implement those standards. \\
- Detection recommendations for text content.

NIST must then report the framework to the Senate Commerce Committee and House Energy \& Commerce Committee with legislative/administrative recommendations.

\textbf{Enforcement}

Violations are treated as \textbf{unfair or deceptive acts or practices} under Sec.18(a)(1)(B) of the FTC Act (15 U.S.C. 57a(a)(1)(B)). The \textbf{FTC} enforces with full FTC Act powers and penalties. The bills explicitly preserve any other authority the FTC already holds.

\textbf{Status:} Both bills were referred to committee in the 118th Congress (S. 2691 to Senate Commerce; H.R. 6466 to the House Subcommittee on Innovation, Data, and Commerce) and did not advance further.

\fieldlabel{Assessment.} \colorbox{passgreen}{\textcolor{white}{\bfseries PASS}}\par
Hit all four grading-key points: identified the H.R. 6466 / S. 2691 companions, the AI-generated-content disclosure requirements, the developer-and-licensee duty to prevent downstream removal, and FTC enforcement. Detailed and citation-grounded. Minor blemish: a leaked interstitial line ("Let me pull both simultaneously") appears in the answer text.
\vspace{4pt}\hrule\vspace{8pt}
\subsection*{Q4 \quad \colorbox{passgreen}{\textcolor{white}{\bfseries PASS}}}
\fieldlabel{Category.} conceptual retrieval\par
\fieldlabel{Question.} Which bills in the corpus address AI-generated deepfakes?\par
\fieldlabel{Expected (corpus ground truth).}\par
\textasciitilde{}10 bills mention deepfakes, including the No AI FRAUD Act (118-hr-6943), Protecting Consumers from Deceptive AI Act (118-hr-7766 / 119-hr-8479), AI PLAN Act, AI Public Awareness and Education Campaign Act, American Leadership in AI Act, and the AI Scam Prevention Act. Several are companion/reintroduced pairs.

\fieldlabel{Agent answer.}\quad{\small\itshape 44.53s, ttft 44.53s, tools: search\_corpus, search\_corpus, search\_corpus, search\_corpus, search\_corpus, search\_corpus, find\_definitions\_of, get\_bill, get\_bill, get\_bill, get\_bill, get\_section, get\_section, get\_section, get\_bill, get\_bill, get\_section, get\_section, get\_section}\par
Here is a full accounting of the bills in the corpus that address AI-generated deepfakes, organized by approach:

\textbf{1. Watermarking \& Provenance Disclosure (Deepfake Detection)}

\textbf{H.R. 7766, 118th Cong.} --- *Protecting Consumers from Deceptive AI Act* (Eshoo, D-CA) \\
\textbf{H.R. 8479, 119th Cong.} --- *Protecting Consumers from Deceptive AI Act* (Foushee, D-NC) *(119th reintroduction)*

These are companion bills (nearly identical text). Both direct NIST to establish task forces on watermarking, digital fingerprinting, and content provenance metadata for AI-generated audio/visual content. They also impose direct obligations on generative AI app providers to embed machine-readable disclosures acknowledging AI origin, and on large online platforms (\textgreater{}=\$50M revenue or \textgreater{}=25M monthly users) to display --- and not remove --- those disclosures. Enforcement falls to the FTC as an unfair/deceptive trade practice. *See* Sec. 3 of H.R. 7766, 118th Cong. and Sec. 3 of H.R. 8479, 119th Cong.

\textbf{2. Deepfakes as Financial Crime / National Security Risk}

\textbf{H.R. 7781, 118th Cong.} --- *AI PLAN Act* (Nunn, R-IA) \\
\textbf{H.R. 2152, 119th Cong.} --- *AI PLAN Act* (Nunn, R-IA) *(119th version; passed committee 52--0 in May 2026)*

Both bills require the Secretaries of Treasury, Homeland Security, and Commerce to submit annual reports to Congress on AI-enabled financial crime risks. Deepfakes are explicitly listed as a risk factor to be analyzed, alongside voice cloning, synthetic identities, foreign election interference, and market disruption. *See* Sec. 2 of H.R. 7781, 118th Cong. and Sec. 2 of H.R. 2152, 119th Cong.

\textbf{3. AI-Generated Likeness / Voice Rights (Deepfake Identity Fraud)}

\textbf{H.R. 6943, 118th Cong.} --- *No AI FRAUD Act* (Salazar, R-FL)

Creates federal property rights in an individual's likeness and voice, including after death. Prohibits distributing "personalized cloning services" or publishing unauthorized "digital depictions" or "digital voice replicas" without consent. Penalties: \$50,000/violation for cloning services; \$5,000/violation for unauthorized depictions/voice replicas; plus actual damages, profits, punitive damages, and attorneys' fees. Sexual/intimate images are classified as per se harm. *See* Sec. 3 of H.R. 6943, 118th Cong.

\textbf{4. Public Awareness (Deepfakes as a Category of Concern)}

\textbf{S. 4596, 118th Cong.} / \textbf{S. 1699, 119th Cong.} / \textbf{H.R. 7151, 119th Cong.} --- *Artificial Intelligence Public Awareness and Education Campaign Act* (Young, R-IN)

Directs Commerce to run a public education campaign on AI risks, explicitly including "media commonly referred to as deepfakes." Also targets outreach to populations vulnerable to AI-enabled fraud. *See* Sec. 2 of S. 4596, 118th Cong.

\textbf{Summary Table}

| Bill | Approach | Key Mechanism | \\
|---|---|---| \\
| H.R. 7766 (118th) / H.R. 8479 (119th) | Detection \& disclosure | Watermarks, provenance metadata, FTC enforcement | \\
| H.R. 7781 (118th) / H.R. 2152 (119th) | Risk reporting | Annual agency reports to Congress on deepfake-enabled financial crime | \\
| H.R. 6943 (118th) | IP / liability | Federal right of likeness/voice; private right of action | \\
| S. 4596 (118th) / S. 1699 \& H.R. 7151 (119th) | Awareness | Public education campaign explicitly naming deepfakes |

No bill in the corpus directly criminalizes deepfakes at the federal level or creates a standalone civil cause of action specifically for political/electoral deepfakes --- those angles are not represented.

\fieldlabel{Assessment.} \colorbox{passgreen}{\textcolor{white}{\bfseries PASS}}\par
Returned a well-organized set of 8 deepfake bills grouped by approach, named the No AI FRAUD Act and the Protecting Consumers from Deceptive AI Act, and explicitly flagged companion/reintroduction pairs. It missed 2 of the 10 ground-truth bills (American Leadership in AI Act, AI Scam Prevention Act) but 8/10 is well within a defensible set.
\vspace{4pt}\hrule\vspace{8pt}
\subsection*{Q5 \quad \colorbox{failred}{\textcolor{white}{\bfseries FAIL}}}
\fieldlabel{Category.} conceptual retrieval (vocabulary gap)\par
\fieldlabel{Question.} Which bills restrict exports of advanced AI chips or semiconductors to foreign adversaries?\par
\fieldlabel{Expected (corpus ground truth).}\par
H.R. 4683 (CLOUD AI Act --- blocks cloud/remote access to advanced ICs for China/Macau entities) and H.R. 5885 (GAIN AI Act --- advanced-IC export license requirement). S. 321 (Decoupling America's AI Capabilities from China Act) bars import/export of AI tech/IP. H.R. 6996 'Full AI Stack Export Promotion Act' is the opposite --- it promotes exports --- and must not be miscounted.

\fieldlabel{Agent answer.}\quad{\small\itshape 23.86s, ttft 23.86s, tools: search\_corpus, search\_corpus, search\_corpus, search\_corpus, get\_bill, get\_bill, search\_corpus, search\_corpus, search\_corpus, search\_corpus}\par
\textbf{The corpus contains no standalone bills specifically restricting AI chip/semiconductor exports to foreign adversaries.} Every search consistently surfaces only two bills, both large omnibus National Defense Authorization Acts. Here's what the corpus shows:

\textbf{Bills with relevant provisions}

\textbf{1. National Defense Authorization Act for FY2024 --- H.R. 2670, 118th Cong.} \\
- Enacted as \textbf{Public Law No. 118-31} (Dec. 22, 2023) \\
- Matches on AI + export controls + foreign adversary, but as one of hundreds of provisions inside a sprawling NDAA. It does not appear to be a dedicated AI chip export bill.

\textbf{2. Servicemember Quality of Life Improvement and National Defense Authorization Act for FY2025 --- H.R. 5009, 118th Cong.} \\
- Enacted as \textbf{Public Law No. 118-159} (Dec. 23, 2024) \\
- Matches on AI + hardware + China + Russia + national security, again embedded within a broad NDAA.

\textbf{Why the gaps exist}

This corpus covers \textbf{191 bills from the 118th and 119th Congresses that primarily or substantially concern AI governance}. Bills whose *primary* focus is export control law (e.g., Export Administration Regulations reform, CHIPS Act implementation rules, Commerce Department licensing frameworks) are \textbf{out of scope} --- even if they touch AI hardware --- because they are classified as trade/national-security legislation, not AI-governance legislation. The Biden and Trump-era AI chip export rules (e.g., the "diffusion rule") are also \textbf{executive/regulatory actions}, which this corpus explicitly excludes.

If you want to drill into the NDAA provisions that matched, I can pull the relevant sections from H.R. 2670 or H.R. 5009 --- just say the word.

\fieldlabel{Assessment.} \colorbox{failred}{\textcolor{white}{\bfseries FAIL}}\par
Concluded the corpus contains "no standalone bills specifically restricting AI chip/semiconductor exports" and surfaced only two NDAAs. The corpus in fact contains at least three on-point bills --- the CLOUD AI Act (H.R. 4683), the GAIN AI Act (H.R. 5885), and S. 321 --- which BM25 keyword search failed to surface under the query phrasing. The agent did not fabricate (it correctly avoided miscounting the export-promotion bill H.R. 6996), but it delivered a wrong answer: it told the user a topic is absent when it is present. The clearest case in the suite for hybrid semantic retrieval.
\vspace{4pt}\hrule\vspace{8pt}
\subsection*{Q6 \quad \colorbox{partialorange}{\textcolor{white}{\bfseries PARTIAL}}}
\fieldlabel{Category.} aggregate count\par
\fieldlabel{Question.} How many bills in the corpus were introduced in the 119th Congress?\par
\fieldlabel{Expected (corpus ground truth).}\par
91 bills were introduced in the 119th Congress.

\fieldlabel{Agent answer.}\quad{\small\itshape 12.78s, ttft 4.39s, tools: corpus\_coverage, search\_corpus}\par
The corpus coverage report gives bill counts by \textbf{tier} but not broken down by Congress. Let me search to get a precise count for the 119th.The `total` field on that search returns 91 hits for the 119th Congress, but that's against a specific keyword query --- not a full count. The `coverage\_note` confirms the corpus has \textbf{191 bills total} across the 118th and 119th Congresses, but the tools don't expose a per-Congress breakdown directly.

What I can say with confidence from the corpus metadata:

- \textbf{Total corpus:} 191 bills (118th + 119th Congress combined), spanning introductions from January 9, 2023 through May 4, 2026. \\
- \textbf{Tier breakdown (all congresses):} 184 Tier A + 7 Tier B. \\
- The tools do not provide a precise per-Congress split in a single aggregate call --- a filtered search for "artificial intelligence" against congress=119 returns 91 matching bills, but that's a keyword-matched subset, not a definitive count of all 119th Congress bills in the corpus.

For a precise 119th-only count, the corpus coverage tool would need to expose per-Congress bill counts --- which it currently does not in v1. The safe answer: \textbf{the corpus holds 191 bills across both congresses; the exact 119th Congress share is not directly reported by the available tools.}

\fieldlabel{Assessment.} \colorbox{partialorange}{\textcolor{white}{\bfseries PARTIAL}}\par
The true answer is 91, and the agent literally had "91" in hand from a search total --- but distrusted it as a keyword-matched subset and declined to give a definitive count. The caution is defensible and it did not fabricate, but it failed to answer a question with a clean answer. Root cause: no per-Congress count primitive exists.
\vspace{4pt}\hrule\vspace{8pt}
\subsection*{Q7 \quad \colorbox{partialorange}{\textcolor{white}{\bfseries PARTIAL}}}
\fieldlabel{Category.} aggregate fact\par
\fieldlabel{Question.} Which member of Congress sponsored the most bills in this corpus, and how many?\par
\fieldlabel{Expected (corpus ground truth).}\par
Rep. Ted Lieu (D-CA-36) sponsored the most, with 10 bills.

\fieldlabel{Agent answer.}\quad{\small\itshape 26.85s, ttft 26.84s, tools: search\_corpus, search\_corpus, search\_corpus, search\_corpus, search\_corpus, search\_corpus, corpus\_coverage, search\_corpus, search\_corpus}\par
The search tool doesn't support sponsor-aggregation queries and the corpus tools don't include a "list all bills by sponsor" primitive. From the search results returned across all my queries, I can tabulate the sponsors that appeared, but I only surfaced \textasciitilde{}17 distinct bills --- a fraction of the 191 in the corpus. The tools available don't include a sponsor-rollup or "find\_bills\_by\_sponsor" function, so I cannot accurately count every bill per sponsor across all 191 entries.

\textbf{What I can say from the data actually returned:}

- \textbf{Rep. Ted Lieu (D-CA-36)} is the only sponsor who appeared \textbf{more than once} in results: he sponsored at least \textbf{3 bills} in the corpus --- \\
  - *National AI Commission Act* (118th, H.R. 4223) \\
  - *Block Nuclear Launch by Autonomous Artificial Intelligence Act of 2023* (118th, H.R. 2894) \\
  - *American Leadership in AI Act* (119th, H.R. 8516)

However, that count is almost certainly incomplete --- those are only the bills that surfaced in keyword searches. The corpus has 191 bills, and the search tool is not designed to enumerate all bills by a given sponsor.

\textbf{Bottom line:} The tools available here don't include a sponsor-aggregation primitive, so I cannot reliably answer "who has the *most* bills" across all 191 entries without risking an inaccurate count. If this feature is important, a `find\_bills\_by\_sponsor` aggregate function --- analogous to `find\_bills\_defining` --- would be needed to answer it reliably.

\fieldlabel{Assessment.} \colorbox{partialorange}{\textcolor{white}{\bfseries PARTIAL}}\par
Identified the right person (Rep. Ted Lieu) but gave the wrong count ("at least 3," vs. the true 10) and explicitly declined to answer definitively, because no sponsor-aggregation tool exists. Honest about the limitation and no fabrication, but it did not deliver the answer. It usefully suggested a find\_bills\_by\_sponsor primitive --- the correct fix.
\vspace{4pt}\hrule\vspace{8pt}
\subsection*{Q8 \quad \colorbox{passgreen}{\textcolor{white}{\bfseries PASS}}}
\fieldlabel{Category.} definitions aggregate\par
\fieldlabel{Question.} How do bills in the corpus define "artificial intelligence"? Are the definitions consistent across bills?\par
\fieldlabel{Expected (corpus ground truth).}\par
57 bills define 'artificial intelligence'. They are largely consistent: 42 define it by reference, and 35 of those point to 15 U.S.C. 9401 (the National AI Initiative Act of 2020 definition). Only 15 write a direct definition, and those track the same structure. The corpus converges on the 15 U.S.C. 9401 statutory definition.

\fieldlabel{Agent answer.}\quad{\small\itshape 31.13s, ttft 31.13s, tools: find\_definitions\_of}\par
\textbf{57 bills} in the corpus define "artificial intelligence." The short answer on consistency: \textbf{the overwhelming majority converge on a single statutory anchor, but a meaningful minority diverge --- and the few bills that write their own definitions vary considerably.}

\textbf{The dominant pattern: deferral to 15 U.S.C. 9401 (\textasciitilde{}70\% of bills)}

Roughly \textbf{42 of 57 bills} (74\%) define AI by reference to \textbf{Sec. 5002 of the National AI Initiative Act of 2020 (15 U.S.C. 9401)}. The language is nearly identical across all of them --- for example:

\textgreater{} "The term *artificial intelligence* has the meaning given such term in section 5002 of the National Artificial Intelligence Initiative Act of 2020 (15 U.S.C. 9401)."

This is the clear legislative default across both the 118th and 119th Congresses. Bills using this approach include H.R. 6881, H.R. 7381, S. 3732, S. 4487, H.R. 1736, H.R. 5045, S. 3063, S. 3952, and dozens more (see full list below).

\textbf{Secondary reference: 10 U.S.C. note prec. 4061 (\textasciitilde{}5 bills)}

A smaller cluster defers instead to \textbf{Sec. 238(g) of the John S. McCain NDAA for FY 2019} (10 U.S.C. note prec. 4061). This is an older, defense-focused definition. Bills using this anchor:

| Bill | Citation | \\
|------|----------| \\
| S. 2293, 118th | Sec. 2(2) of S. 2293, 118th Cong. | \\
| S. 2346, 118th | Sec. 2(a)(1) of S. 2346, 118th Cong. | \\
| S. 501, 119th | Sec. 2(a)(1) of S. 501, 119th Cong. | \\
| S. 2381, 119th | Sec. 2(1) of S. 2381, 119th Cong. | \\
| S. 3202, 119th | Sec. 2(f)(1) of S. 3202, 119th Cong. |

\textbf{Hybrid definitions (\textasciitilde{}2 bills)}

Two bills layer \textbf{both} statutory sources:

- \textbf{S. 4495, 118th Cong.} (Sec. 2(3)): Adopts 15 U.S.C. 9401 *and* incorporates the NDAA FY2019 techniques (paragraphs 1--5 of Sec. 238(g)). \\
- \textbf{H.R. 4142, 119th Cong.} (Sec. 3(d)(2)): Same dual-reference approach.

\textbf{Standalone direct definitions (\textasciitilde{}7 bills)}

These bills write their own substantive text --- and they differ materially:

| Bill | Definition | \\
|------|-----------| \\
| \textbf{H.R. 7913, 118th} --- Sec. 2(d)(1) | "An automated system designed to perform a task typically associated with human intelligence or cognitive function." --- Very broad, cognitive-function framing. | \\
| \textbf{H.R. 8005, 118th} --- Sec. 2(j)(1) | "A machine-based system that uses data, references, or machine or human-based inputs to (A) create a visual depiction; (B) recognize a pattern; or (C) make a recommendation, prediction, or decision." --- Narrower, output-focused, with an enumerated list. | \\
| \textbf{H.R. 9673, 118th / S. 4838, 118th} --- Sec. 3(1) each | "An engineered system that can, for a given set of human-defined objectives, generate output such as content, predictions, recommendations, or decisions that influence people or environments." --- Identical text in both bills. | \\
| \textbf{H.R. 3210, 119th} --- Sec. 2(c)(2) | Near-identical to H.R. 9673/S. 4838 but with minor word-order variation. | \\
| \textbf{S. 4896, 118th} --- Sec. 2(a)(1) | "A machine-based system that can, for a given set of human-defined objectives, make predictions, recommendations or decisions... [using] machine and human-based inputs to perceive environments, abstract perceptions into models, and use model inference to formulate options." --- Most detailed standalone definition; closest to OECD-style language. | \\
| \textbf{S. 1792, 119th / H.R. 8516, 119th} --- Sec. 2(3) / Sec. V(C)(521)(3) | Five-part enumeration covering systems that learn from data, solve human-like tasks, think/act like humans, use ML techniques, or act rationally --- the most expansive in the corpus. | \\
| \textbf{S. 321, 119th} --- Sec. 2(1) | Cross-references 18 U.S.C. 2741 (a new provision the bill itself creates), tying "artificial intelligence" to its own novel statutory definition of generative AI. |

\textbf{Bottom line}

\textbf{Consistency is high on paper but incomplete in substance.} The 15 U.S.C. 9401 anchor dominates, creating formal uniformity across \textasciitilde{}74\% of bills. But that uniformity is partly illusory --- it just defers disagreement to the underlying NAIIA definition rather than resolving it. The NDAA FY2019 cluster uses a meaningfully different, more enumerated definition. The standalone drafters diverge most: H.R. 7913 relies on a cognitive-function test; H.R. 8005 is output-and-task focused; S. 1792/H.R. 8516 use a five-category list; S. 4896 is procedure-focused. No single Congress-wide definition has been enacted.

\fieldlabel{Assessment.} \colorbox{passgreen}{\textcolor{white}{\bfseries PASS}}\par
Correct and unusually strong: 57 bills define "artificial intelligence," the majority by reference, \textasciitilde{}42 pointing to 15 U.S.C. 9401. The nuance --- formal uniformity that defers, rather than resolves, substantive disagreement --- is better analysis than the expected answer asked for.
\vspace{4pt}\hrule\vspace{8pt}
\subsection*{Q9 \quad \colorbox{passgreen}{\textcolor{white}{\bfseries PASS}}}
\fieldlabel{Category.} amendment graph\par
\fieldlabel{Question.} Which bills in the corpus amend the Communications Act of 1934?\par
\fieldlabel{Expected (corpus ground truth).}\par
H.R. 8939 and H.R. 334 (companion bills amending Sec.227(d)(3) for AI-generated voice), S. 1993 (amends Sec.230 to waive interactive-computer-service immunity for generative-AI claims), H.R. 7786 (AI Fraud Accountability Act, Sec.223), and S. 3495 (AI Scam Prevention Act, Sec.Sec.227 \& 230). Bills that merely mention the Act are not amenders.

\fieldlabel{Agent answer.}\quad{\small\itshape 31.46s, ttft 7.19s, tools: search\_corpus, resolve\_citation, get\_amendments, get\_amendments, get\_amendments, get\_amendments, get\_amendments, get\_amendments}\par
Now let me check the amendments data for the most promising candidates, and also get the amendments for the key bills that explicitly target the Communications Act:Here is the complete picture of AI-governance corpus bills that amend provisions of the Communications Act of 1934 (codified in Title 47). All amendment target texts are flagged `target\_text\_unverified` pending USC ingestion.

\textbf{Bills Amending the Communications Act of 1934}

The Communications Act is codified at Title 47 U.S.C. The relevant provisions targeted are Sec.Sec. 223, 227, and 230.

\textbf{1. \textbf{H.R. 7786 / S. 3982 --- AI Fraud Accountability Act} (119th Cong., 2026)} \\
\textbf{Target: 47 U.S.C. Sec. 223 (Obscene or harassing telephone calls)}

Both the House and Senate companion bills insert a new subsection into Sec. 223 criminalizing the use of \textbf{AI-generated digital impersonations to commit fraud} over interstate or foreign communications. Key provisions: \\
- Defines "digital impersonation" (AI-generated visual/audio indistinguishable from a real person) and "identifiable individual." \\
- Makes it unlawful to falsely pose as a real or imaginary individual via a digital impersonation with intent to defraud. \\
- Penalties: up to 3 years imprisonment, fines, and criminal forfeiture. \\
- Extraterritorial federal jurisdiction.

*Cited amendments:* Sec. 2(a)(2) of H.R. 7786, 119th Cong. -\textgreater{} 47 U.S.C. 223; Sec. 2(a)(2) of S. 3982, 119th Cong. (target statute unresolved in corpus for S. 3982, but text is identical).

\textbf{2. \textbf{S. 3495 --- Artificial Intelligence Scam Prevention Act} (119th Cong., 2025)} \\
\textbf{Target: 47 U.S.C. Sec. 227 (Telephone Consumer Protection Act / TCPA) and 47 U.S.C. Sec. 230}

Three operations amend Sec. 227: \\
- \textbf{Sec. 4(b)(1)(A)(ii):} Expands the definition of "artificial or prerecorded voice" to include any AI-generated technological representation of speech that appears to authentically depict an individual who did not engage in that speech. \\
- \textbf{Sec. 4(b)(1)(A)(iv):} Adds definitions of "text message" and "text messaging service" to Sec. 227. \\
- \textbf{Sec. 4(b)(1)(C):} Adds a new subsection requiring disclosure at the start of any autodialed call or text that uses AI to emulate a human being.

One operation amends Sec. 230: \\
- \textbf{Sec. 4(b)(2):} Adds a clarification that Sec. 230 does not limit application of the Telemarketing and Consumer Fraud and Abuse Prevention Act.

*(The bill also amends 15 U.S.C. Sec.Sec. 6102 and 6106, outside the Communications Act.)*

\textbf{3. \textbf{S. 1993 --- No Section 230 Immunity for AI Act} (118th Cong., 2023)} \\
\textbf{Target: 47 U.S.C. Sec. 230 (Section 230 liability shield)}

Two operations: \\
- \textbf{Sec. 1(1):} Adds a new paragraph stating that Sec. 230 immunity (except for Sec. 230(c)(2)(A)) does \textbf{not} apply to civil or criminal claims where the conduct involves the "use or provision of generative artificial intelligence" by an interactive computer service. \\
- \textbf{Sec. 1(2):} Adds a definition of "generative artificial intelligence" to Sec. 230's definitions --- "an AI system capable of generating novel text, video, images, audio, and other media based on prompts."

\textbf{4. \textbf{H.R. 8939 / H.R. 334 --- AI voice standards bills} (118th/119th Cong.)} \\
Both bills are \textbf{titled} as amendments to the Communications Act (establishing technical/procedural standards for AI-generated prerecorded voice systems), but no amendment operations were extracted from their XML --- the corpus returned 0 operations for each. The substantive content is in the bill text but was not parsed into discrete amendment records.

\textbf{Summary table:}

| Bill | Congress | Target Statute | Subject | \\
|---|---|---|---| \\
| H.R. 7786 | 119th | 47 U.S.C. Sec. 223 | AI digital impersonation fraud | \\
| S. 3982 | 119th | 47 U.S.C. Sec. 223 | AI digital impersonation fraud (companion) | \\
| S. 3495 | 119th | 47 U.S.C. Sec.Sec. 227, 230 | AI robocall/text disclosure; TCPA expansion | \\
| S. 1993 | 118th | 47 U.S.C. Sec. 230 | Strip Sec. 230 immunity for generative AI | \\
| H.R. 8939 | 118th | Communications Act | AI prerecorded voice standards (no parsed ops) | \\
| H.R. 334 | 119th | Communications Act | AI prerecorded voice standards (no parsed ops) |

\fieldlabel{Assessment.} \colorbox{passgreen}{\textcolor{white}{\bfseries PASS}}\par
Identified the Communications Act amenders correctly, including the three ground-truth bills (H.R. 334, H.R. 8939, S. 1993), the correct target sections (Sec.227 voice systems, Sec.230 immunity, Sec.223 fraud), and S. 3495. It honestly flagged that H.R. 8939/334 returned zero parsed amendment operations despite being titled as amendments --- faithful about an extraction-layer gap rather than papering over it.
\vspace{4pt}\hrule\vspace{8pt}
\subsection*{Q10 \quad \colorbox{passgreen}{\textcolor{white}{\bfseries PASS}}}
\fieldlabel{Category.} definitions (honest-null probe)\par
\fieldlabel{Question.} Which bills in the corpus, if any, define the term "frontier model"?\par
\fieldlabel{Expected (corpus ground truth).}\par
No bill in the corpus defines 'frontier model'. The definition index and full-text search both return zero matches. (S. 5616 uses the distinct term 'covered frontier artificial intelligence model'.)

\fieldlabel{Agent answer.}\quad{\small\itshape 6.26s, ttft 6.26s, tools: find\_bills\_defining}\par
No bills in the corpus formally define "frontier model" (or the close variants "frontier AI model" / "frontier AI system"). The term appears in policy discourse but has not been codified in a definition section in any of the 191 AI-governance bills covered by the corpus.

If you want, I can search for bills that *use* the phrase "frontier model" in their text, even without a formal definition --- just say the word.

\fieldlabel{Assessment.} \colorbox{passgreen}{\textcolor{white}{\bfseries PASS}}\par
Correctly answered that no bill in the corpus defines "frontier model," did not fabricate a definition, and offered to search for usages instead. Clean honest-null behavior.
\vspace{4pt}\hrule\vspace{8pt}
\subsection*{Q11 \quad \colorbox{passgreen}{\textcolor{white}{\bfseries PASS}}}
\fieldlabel{Category.} single-bill substantive\par
\fieldlabel{Question.} What does the Generative AI Terrorism Risk Assessment Act require?\par
\fieldlabel{Expected (corpus ground truth).}\par
H.R. 1736 (119th Congress) requires DHS, with the DNI, to deliver an annual assessment to Congress for five years of terrorism threats from terrorist use of generative AI. It also directs DHS to review/disseminate fusion-center intelligence and obliges agencies to share information with DHS. It is a reporting mandate, not an AI restriction.

\fieldlabel{Agent answer.}\quad{\small\itshape 18.79s, ttft 18.79s, tools: search\_corpus, get\_bill, get\_section, get\_section}\par
\textbf{H.R. 1736, 119th Cong. --- Generative AI Terrorism Risk Assessment Act} \\
Sponsor: Rep. Pfluger (R-TX). Passed the House (Nov. 19, 2025); referred to Senate Homeland Security \& Governmental Affairs Committee (Nov. 20, 2025).

\textbf{Core Requirement: Annual Threat Assessments}

The Secretary of Homeland Security, \textbf{in consultation with the Director of National Intelligence}, must submit an assessment of generative AI terrorism threats to Congress \textbf{within one year of enactment and annually for five years thereafter}. (Sec. 3 of H.R. 1736, 119th Cong.)

Each assessment must include:

1. \textbf{Incident analysis} --- a review of the preceding calendar year's incidents in which a foreign terrorist organization or individual *used or attempted to use* generative AI to: \\
   - Spread violent extremist messaging and facilitate radicalization/recruitment; or \\
   - Enhance development or deployment of \textbf{chemical, biological, radiological, or nuclear (CBRN) weapons}.

2. \textbf{Countermeasure recommendations} --- appropriate measures to counter the threats identified.

\textbf{Procedural Requirements}

- \textbf{Form:} Each assessment must be submitted in \textbf{unclassified form} (a classified annex is permitted). The unclassified portion must be \textbf{posted publicly} on DHS's website and cannot contain FOUO-designated or otherwise non-public information. \\
- \textbf{Briefing:} Within \textbf{30 days} of each submission, DHS must brief the appropriate congressional committees. Other relevant federal agency heads may be required to join. \\
- \textbf{Privacy/civil liberties:} DHS must coordinate assessments to ensure compliance with applicable law and protections for privacy, civil rights, and civil liberties.

\textbf{Information Sharing}

DHS must \textbf{review fusion center intelligence} on generative AI terrorism threats and incorporate it into DHS's own assessments, and \textbf{disseminate} DHS information back to state and major urban area fusion centers. The following entities are required to share relevant information with DHS: \\
- Office of the Director of National Intelligence \\
- FBI \\
- Members of the intelligence community \\
- Any other relevant federal agency DHS deems necessary

\textbf{Sense of Congress (Sec. 2)}

Congress declares that the use of generative AI by foreign terrorist organizations is a \textbf{national security threat that is not well understood}, and that DHS --- in consultation with ODNI --- must take steps to recognize, assess, and address it.

\textbf{Key Definitions (Sec. 3(c))}

- \textbf{Generative AI:} "the class of artificial intelligence models that emulate the structure and characteristics of input data in order to generate derived synthetic content, including images, videos, audio, text, and other digital content." \\
- \textbf{Artificial intelligence:} incorporates the definition from 15 U.S.C. 9401 (National AI Initiative Act of 2020). \\
- \textbf{Foreign terrorist organization:} any entity designated under 8 U.S.C. 1189.

\fieldlabel{Assessment.} \colorbox{passgreen}{\textcolor{white}{\bfseries PASS}}\par
Accurate and detailed: DHS in consultation with the DNI must deliver assessments within one year and annually for five years, with the fusion-center and interagency information-sharing requirements correctly described, and the bill correctly framed as a reporting mandate. Added correct supporting detail (CBRN, unclassified-form requirement).
\vspace{4pt}\hrule\vspace{8pt}
\subsection*{Q12 \quad \colorbox{passgreen}{\textcolor{white}{\bfseries PASS}}}
\fieldlabel{Category.} long comparison\par
\fieldlabel{Question.} Compare the AI disclosure obligations in the AI Labeling Act versus the Generative AI Copyright Disclosure Act of 2024. What does each require, and who is the target of each?\par
\fieldlabel{Expected (corpus ground truth).}\par
The AI Labeling Act requires outward, consumer-facing labeling --- generative AI systems must mark output as AI-generated, developers/licensees enforce non-removal, FTC oversight. The Generative AI Copyright Disclosure Act of 2024 (H.R. 7913) requires upstream, regulator-facing disclosure --- whoever builds or significantly alters an AI training dataset must file a notice with the Copyright Office summarizing copyrighted works used, \textgreater{}=30 days before public release.

\fieldlabel{Agent answer.}\quad{\small\itshape 51.65s, ttft 5.65s, tools: search\_corpus, search\_corpus, search\_corpus, get\_bill, search\_corpus, get\_section, search\_corpus, search\_corpus, search\_corpus, search\_corpus, get\_bill, get\_bill, get\_section, get\_section}\par
The "AI Labeling Act" didn't surface directly. Let me search more specifically and pull both bills' full details simultaneously.The "AI Labeling Act" isn't surfacing by name. Let me try a more targeted search and pull the copyright bill's section text simultaneously.There are two versions of the AI Labeling Act of 2023 --- a House bill (H.R. 6466) and a Senate companion (S. 2691). Let me pull both bills' sections now.I now have the full text of all three bills. Here is the complete, citation-grounded comparison.

\textbf{AI Labeling Act of 2023 vs. Generative AI Copyright Disclosure Act of 2024}

\textbf{Note:} The AI Labeling Act exists as identical companion bills --- H.R. 6466 (House, Rep. Kean, introduced Nov. 2023) and S. 2691 (Senate, Sen. Schatz, introduced Jul. 2023). Their substantive disclosure sections are word-for-word identical. The comparison below treats them as one law.

\textbf{AI Labeling Act of 2023 (H.R. 6466 / S. 2691)}

\textbf{Who is targeted:} Developers and operators of generative AI systems that produce content *in commerce*, plus third-party licensees who deploy those systems.

\textbf{What it requires:}

1. \textbf{Content-level labels on outputs} --- Every generative AI system producing image, video, audio, or multimedia content in interstate commerce must attach a \textbf{clear and conspicuous disclosure} to each output identifying it as AI-generated. The bill specifies detailed criteria: the disclosure must match the medium (visual or audible), be prominent, use plain language, and be non-avoidable. *(Sec. 2(a)(1), Sec. 2(c)(3))*

2. \textbf{Metadata embedding} --- For image/video/audio/multimedia outputs, the output's metadata must include: (a) identification as AI-generated, (b) the identity of the tool used, and (c) the date and time of creation. *(Sec. 2(a)(1)(B)(ii))*

3. \textbf{Text content disclosure} --- AI systems producing text (including chatbots) must include a clear and conspicuous disclosure identifying the content as AI-generated. *(Sec. 2(a)(2))*

4. \textbf{Permanence obligation} --- All disclosures must be, to the extent technically feasible, \textbf{permanent or unable to be easily removed} by subsequent users. *(Sec. 2(a)(1)(B)(iii), 2(a)(2))*

5. \textbf{Developer contractual obligations} --- Developers must implement reasonable procedures to prevent downstream disclosure-stripping: requiring contracts and certifications from end users/licensees not to remove labels, and terminating access upon known violations. *(Sec. 2(a)(3)(A))*

6. \textbf{Third-party licensee obligations} --- Licensees face the same downstream-protection duties toward their own end users. *(Sec. 2(a)(3)(B))*

7. \textbf{Enforcement:} Violations are treated as \textbf{FTC Act unfair or deceptive acts}, enforced by the FTC. *(Sec. 2(a)(4))*

8. \textbf{NIST working group:} NIST must convene a working group within 90 days to develop technical detection standards and report to Congress within one year. *(Sec. 2(b))*

\textgreater{} *Citations: Sec. 2 of H.R. 6466, 118th Cong. / Sec. 2 of S. 2691, 118th Cong.*

\textbf{Generative AI Copyright Disclosure Act of 2024 (H.R. 7913)}

\textbf{Who is targeted:} Any \textbf{person who creates or significantly alters a training dataset} used to build a generative AI system --- squarely aimed at AI developers and trainers, not at content consumers or platforms.

\textbf{What it requires:}

1. \textbf{Pre-launch notice to the Copyright Office} --- Before (or within 30 days after enactment for existing systems) a generative AI system is made available to consumers, the person who created or significantly altered its training dataset must file a notice with the \textbf{Register of Copyrights} containing: \\
   - A \textbf{sufficiently detailed summary of any copyrighted works} used in training (or in the significant alteration); and \\
   - The \textbf{URL of the training dataset}, if it is publicly available on the internet. *(Sec. 2(a)(1))*

2. \textbf{Timing:} Notices for new systems due \textbf{30 days before consumer launch}; notices for pre-enactment systems due \textbf{30 days after the Act takes effect}. *(Sec. 2(a)(2))*

3. \textbf{Public database:} The Register of Copyrights must establish and maintain a \textbf{publicly accessible online database} of all submitted notices. *(Sec. 2(c))*

4. \textbf{Enforcement:} Non-compliance triggers a \textbf{civil penalty of not less than \$5,000} per violation. The Register must issue implementing regulations within 180 days. *(Sec. 2(b))*

\textgreater{} *Citation: Sec. 2 of H.R. 7913, 118th Cong.*

\textbf{Side-by-Side Summary}

| Dimension | AI Labeling Act (H.R. 6466 / S. 2691) | Gen. AI Copyright Disclosure Act (H.R. 7913) | \\
|---|---|---| \\
| \textbf{Target} | Developers, operators, licensees deploying AI systems | Persons creating/altering AI training datasets | \\
| \textbf{What is disclosed} | The *output* --- labeling content as AI-generated | The *inputs* --- copyrighted works used in training | \\
| \textbf{To whom} | End users / consumers (on the content itself) | Register of Copyrights (regulatory filing) | \\
| \textbf{Modalities covered} | Image, video, audio, multimedia, text | Any training dataset for any generative AI system | \\
| \textbf{Mechanism} | On-content label + metadata embedding | Notice filing -\textgreater{} public Copyright Office database | \\
| \textbf{Timing} | At point of output generation, continuously | Once, pre-launch (or within 30 days of enactment) | \\
| \textbf{Enforcer} | FTC (unfair/deceptive acts framework) | Register of Copyrights (civil penalties \textgreater{}= \$5,000) | \\
| \textbf{Downstream duty} | Yes --- developers must contractually bind licensees | No analog |

\textbf{Key Analytical Distinction}

The two bills operate at \textbf{opposite ends of the AI pipeline}. The AI Labeling Act is a \textbf{consumer-facing transparency law}: it governs what people *see* when they encounter AI-generated content. The Generative AI Copyright Disclosure Act is a \textbf{creator-facing IP accountability law}: it governs what rightsholders and the public can *learn* about what copyrighted works trained a system, before that system ever reaches a consumer. Neither bill cross-references or builds on the other.

\fieldlabel{Assessment.} \colorbox{passgreen}{\textcolor{white}{\bfseries PASS}}\par
Strong long-form comparison: the AI Labeling Act as consumer-facing output labeling versus H.R. 7913 as a regulator-facing filing with the Register of Copyrights about training-data copyrighted works, with a clean side-by-side table and the correct "opposite ends of the AI pipeline" framing. Minor blemish: leaked interstitial search narration at the start of the answer.
\vspace{4pt}\hrule\vspace{8pt}
\subsection*{Q13 \quad \colorbox{passgreen}{\textcolor{white}{\bfseries PASS}}}
\fieldlabel{Category.} session memory (multi-turn)\par
\fieldlabel{Question (turn 1).} What is the AI for National Security Act?\par
\fieldlabel{Question (turn 2).} Who introduced it, and which Congress was it in?\par
\fieldlabel{Expected (corpus ground truth).}\par
Turn 1: H.R. 1718 amends the NDAA FY2022 to modify DoD enterprise-wide cyber-data procurement, providing for AI-based endpoint security without constant internet connectivity. Turn 2 (memory): must resolve 'it' to that same bill --- sponsor Rep. Jay Obernolte (R-CA), 118th Congress (introduced March 22, 2023).

\fieldlabel{Agent answer (turn 1).}\quad{\small\itshape 15.65s, ttft 15.64s, tools: search\_corpus, get\_bill, get\_section}\par
\textbf{AI for National Security Act --- H.R. 1718, 118th Congress}

- \textbf{Sponsor:} Rep. Jay Obernolte (R-CA-23), with bipartisan cosponsors Reps. Jimmy Panetta (D-CA-19) and Patrick Ryan (D-NY-18) \\
- \textbf{Introduced:} March 22, 2023 \\
- \textbf{Status:} Referred to the House Committee on Armed Services; no further action on record.

\textbf{What it does:}

The bill is narrow and targeted --- a single operative section. It amends \textbf{Section 1521 of the National Defense Authorization Act for FY2022 (P.L. 117-81; 10 U.S.C. 2224 note)}, which governs the Department of Defense's enterprise-wide procurement of cyber data products and services. Specifically, per \textbf{Sec. 2 of H.R. 1718, 118th Cong.}, it makes two changes:

1. \textbf{AI-based endpoint security:} Inserts language explicitly authorizing the use of *"artificial intelligence-based endpoint security that prevents cyber attacks and does not require constant internet connectivity to function"* as part of DoD's enterprise cyber procurement. \\
2. \textbf{Software supply chain security:} Adds language directing DoD to act in its best interests, *"including by enhancing the security of the software supply chain of the Department."*

\textbf{Bottom line:} This is a surgical amendment to existing DoD cyber procurement law, not a sweeping AI governance bill. Its core contribution is formally blessing AI-based endpoint security tools --- particularly offline-capable ones --- as a legitimate component of DoD's enterprise cybersecurity procurement strategy.

\fieldlabel{Agent answer (turn 2).}\quad{\small\itshape 2.46s, ttft 2.46s, tools: \textit{none}}\par
H.R. 1718 was introduced in the \textbf{118th Congress} by \textbf{Rep. Jay Obernolte (R-CA-23)}, with Reps. Jimmy Panetta (D-CA-19) and Patrick Ryan (D-NY-18) as cosponsors. I covered this in my previous answer --- no new lookup needed.

\fieldlabel{Assessment.} \colorbox{passgreen}{\textcolor{white}{\bfseries PASS}}\par
Exemplary session memory. Turn 1 correctly identified H.R. 1718 and its DoD cyber-procurement amendment. Turn 2 resolved "it" to the same bill, answered sponsor (Rep. Jay Obernolte) and Congress (118th) from context, made zero tool calls, and explicitly noted no new lookup was needed.
\vspace{4pt}\hrule\vspace{8pt}
\subsection*{Q14 \quad \colorbox{partialorange}{\textcolor{white}{\bfseries PARTIAL}}}
\fieldlabel{Category.} session memory (multi-turn)\par
\fieldlabel{Question (turn 1).} Which bills in the corpus mention facial recognition?\par
\fieldlabel{Question (turn 2).} Summarize the first one for me.\par
\fieldlabel{Expected (corpus ground truth).}\par
Turn 1: three bills --- ASSESS AI Act (S. 1356, 118th), AI Foundation Model Transparency Act of 2023 (H.R. 6881, 118th), LIFE with AI Act (S. 3063, 119th). Turn 2 (memory): must summarize one of those three (ideally S. 1356, earliest) without re-running the search from scratch or drifting to an unrelated bill.

\fieldlabel{Agent answer (turn 1).}\quad{\small\itshape 11.61s, ttft 11.61s, tools: search\_corpus, search\_corpus}\par
Three bills in the corpus meaningfully mention facial recognition (the biometric searches returned lower-tier NDAA hits with only tangential relevance):

\textbf{Primary hits (Tier A)}

1. \textbf{LIFE with AI Act --- S. 3063, 119th Cong.} *(Sen. Cassidy, R-LA; introduced 2025-10-28)* \\
Highest relevance score. Addresses facial recognition as part of a broader AI lifecycle framework.

2. \textbf{ASSESS AI Act --- S. 1356, 118th Cong.} *(Sen. Bennet, D-CO; introduced 2023-04-27)* \\
Directs a presidential task force to assess privacy, civil rights, and civil liberties implications of AI --- facial recognition is squarely within its mandate. The summary explicitly calls out associated data and civil liberties gaps.

3. \textbf{AI Foundation Model Transparency Act of 2023 --- H.R. 6881, 118th Cong.} *(Rep. Beyer, D-VA; introduced 2023-12-22)* \\
Facial recognition appears in the context of transparency disclosures for foundation models capable of biometric tasks.

\textbf{Lower-tier mentions}

- \textbf{NDAA FY2024 (H.R. 2670)} and \textbf{NDAA FY2025 (H.R. 5009)} matched on "face" + "identification" keywords but are Tier B omnibus defense bills where facial recognition is incidental, not a primary focus.

Want me to drill into the specific sections of any of these bills that address facial recognition?

\fieldlabel{Agent answer (turn 2).}\quad{\small\itshape 34.15s, ttft 4.59s, tools: search\_corpus, get\_bill, get\_section, get\_section, get\_section, get\_section, get\_section, get\_section, get\_section, get\_section, get\_section, get\_section}\par
I need to correct myself --- the previous response was partially fabricated. I should have used the tools first before summarizing. Let me pull the actual bill now.\#\# LIFE with AI Act --- S. 3063, 119th Cong. \\
\textbf{Sponsor:} Sen. Cassidy (R-LA) | \textbf{Introduced:} Oct. 28, 2025 | \textbf{Status:} Referred to Senate HELP Committee

Full title: *Learning Innovation and Family Empowerment with AI Act.* A K--12-focused bill amending FERPA and the Elementary and Secondary Education Act (ESEA) to govern AI and ed-tech use in schools. Ten sections:

\textbf{Definitions (Sec. 2)} \\
- "Artificial intelligence" is imported by reference from 15 U.S.C. 9401 (the NAIA Act definition). \\
- Introduces two new terms: \textbf{"educational technology"} (physical or virtual products/services supporting learning, excluding curriculum) and \textbf{"personalized learning"} (AI-tailored instruction, pacing, and content for individual students). \\
- \textbf{"Instant verification technology"}: real-time parental notification/consent systems operable across smartphones, tablets, and computers.

\textbf{Parental Consent \& Notification (Sec. 3)} \\
Creates a \textbf{"Golden Seal of Excellence in Student Data Privacy"} --- a voluntary certification the Secretary of Education awards to schools that implement instant verification technology for parental consent. Schools must not have violated FERPA in the prior 5 years and must have run the system for at least one academic year. The Seal expires after 5 years.

\textbf{Directory Information Opt-Out (Sec. 4)} \\
Amends FERPA (20 U.S.C. 1232g(a)(5)) to require that any opt-out form be: (1) accessible year-round on multiple digital devices, (2) completable in under 5 minutes, and (3) available on the school's website.

\textbf{Facial Recognition --- the Core AI Restriction (Sec. 5)} \\
Amends FERPA by adding two new subsections: \\
- \textbf{Sec.444(k):} Schools cannot use student photos to train facial recognition AI \textbf{without prior parental consent}, and cannot contract with yearbook companies that use facial recognition unless those companies disclose that use and obtain parental consent. \\
- \textbf{Sec.444(l):} Schools cannot contract with yearbook companies that \textbf{sell student data} collected during yearbook production. \\
*(Sec. 5 of S. 3063, 119th Cong.)*

\textbf{FERPA "Education Record" Redefinition (Sec. 6)} \\
Expands the definition of "education records" to cover \textbf{any data} related to a student (academic performance, attendance, health, discipline) maintained by a school or any entity acting on its behalf --- broadening FERPA's reach to third-party ed-tech vendors.

\textbf{Third-Party Ed-Tech Contracts (Sec. 7)} \\
Schools must publicly post ed-tech contracts for \textbf{at least 2 weeks} before signing. Contracts must require vendor certification of FERPA compliance. The Secretary maintains a \textbf{public list of non-compliant vendors} for 5 years, with an appeals process.

\textbf{Privacy Technical Assistance Center (Sec. 8)} \\
Directs the Secretary to establish a \textbf{Privacy Technical Assistance Center} within 6 months of enactment to help schools comply with FERPA and the Protection of Pupil Rights Amendment. The Center can approve voluntary \textbf{safe harbor programs} run by independent organizations; participants in good standing get a presumption of compliance.

\textbf{AI Integration in Classrooms (Sec. 9)} \\
Directs the Secretary (with the Institute of Education Sciences) to develop AI training resources for K--12 teachers. Amends ESEA (20 U.S.C. 6613) to add AI literacy --- including harms of copyright piracy and improper student AI use --- to teacher professional development requirements.

\textbf{Personalized Learning R\&D (Sec. 10)} \\
Directs the Secretary to prioritize \textbf{AI-in-education} proposals under the Small Business Innovation Research (SBIR) program, favoring projects that enhance personalized learning, serve students with special education needs, and preserve critical thinking skills.

\textbf{Bottom line:} This is primarily a student-privacy bill that grafts AI-specific rules onto FERPA. The facial recognition ban (Sec. 5) is its most concrete AI restriction; the rest of the bill builds out a consent and vendor-accountability infrastructure around ed-tech broadly.

\fieldlabel{Assessment.} \colorbox{partialorange}{\textcolor{white}{\bfseries PARTIAL}}\par
Memory mechanics passed: turn 1 listed the three facial-recognition bills, and turn 2 summarized "the first one" (S. 3063, which the agent had listed first), staying anchored to turn 1's result set. But turn 2 opened with the agent's own admission --- "the previous response was partially fabricated" --- because turn 1 characterized all three bills' provisions from search snippets without pulling their text. The self-correction is healthy, but turn 1 should not have described un-retrieved bills in the first place.
\vspace{4pt}\hrule\vspace{8pt}
\subsection*{Q15 \quad \colorbox{passgreen}{\textcolor{white}{\bfseries PASS}}}
\fieldlabel{Category.} out-of-corpus (state law)\par
\fieldlabel{Question.} What does the California AI Transparency Act require?\par
\fieldlabel{Expected (corpus ground truth).}\par
Out of scope. The corpus is U.S. federal legislation only; the California AI Transparency Act (SB 942) is a California state statute and is not in the corpus. The agent should decline to answer from the corpus rather than fabricate provisions.

\fieldlabel{Agent answer.}\quad{\small\itshape 7.28s, ttft 7.28s, tools: search\_corpus}\par
The "California AI Transparency Act" is \textbf{not in this corpus}. The corpus covers 191 federal AI-governance bills from the 118th and 119th U.S. Congress --- the two results returned are unrelated NDAA bills that merely contain those words scattered across their text.

The California AI Transparency Act is \textbf{state legislation}, and this database is scoped exclusively to \textbf{federal legislation}. State bills, regulatory actions, and executive orders are explicitly out of scope.

For the California AI Transparency Act (SB 942, signed into law in 2024), you'd need a source covering California state legislation --- e.g., the California Legislative Information site (leginfo.legislature.ca.gov) or a state-law database.

\fieldlabel{Assessment.} \colorbox{passgreen}{\textcolor{white}{\bfseries PASS}}\par
Exemplary out-of-scope handling. Correctly identified the California AI Transparency Act as state law (SB 942), stated it is not in the federal corpus, declined to invent provisions, and pointed the user to the right external source.
\vspace{4pt}\hrule\vspace{8pt}
\subsection*{Q16 \quad \colorbox{passgreen}{\textcolor{white}{\bfseries PASS}}}
\fieldlabel{Category.} out-of-corpus (foreign law)\par
\fieldlabel{Question.} Summarize the EU AI Act's risk tiers.\par
\fieldlabel{Expected (corpus ground truth).}\par
Out of scope. The corpus holds only U.S. federal legislation; the EU AI Act is European Union law. The agent should flag it as out of scope rather than recite the risk tiers as if corpus-grounded.

\fieldlabel{Agent answer.}\quad{\small\itshape 4.0s, ttft 3.99s, tools: \textit{none}}\par
The EU AI Act is outside this corpus. The database covers \textbf{US federal legislation only} --- 191 bills from the 118th and 119th Congress (2023--present) that primarily concern AI, ML, generative AI, frontier models, automated decision systems, or facial recognition. EU regulations, including the EU AI Act, are not included and cannot be summarized from these tools.

If you're interested in \textbf{US legislative analogues} --- bills that establish risk-based frameworks for AI (e.g., high-risk AI systems, prohibited uses, tiered obligations) --- I can search for those. Just say the word.

\fieldlabel{Assessment.} \colorbox{passgreen}{\textcolor{white}{\bfseries PASS}}\par
Correctly flagged the EU AI Act as foreign law outside the US-federal corpus and declined to recite its risk tiers from training data, instead offering to find US analogues. Exactly the right behavior for the anti-hallucination thesis.
\vspace{4pt}\hrule\vspace{8pt}
\subsection*{Q17 \quad \colorbox{passgreen}{\textcolor{white}{\bfseries PASS}}}
\fieldlabel{Category.} out-of-corpus (nonexistent bill)\par
\fieldlabel{Question.} What does H.R. 9999 from the 119th Congress say about AI?\par
\fieldlabel{Expected (corpus ground truth).}\par
There is no bill 119-hr-9999 in the corpus (the highest 119th-Congress H.R. number present is 8664). The agent should say it is not in the corpus and not fabricate content.

\fieldlabel{Agent answer.}\quad{\small\itshape 10.13s, ttft 2.19s, tools: search\_corpus, get\_bill, get\_bill}\par
Let me look that up for you.\textbf{H.R. 9999 (119th Congress) is not in the corpus.} There are two possible explanations:

1. \textbf{The bill doesn't exist (yet).} Bill numbers this high are rarely assigned early in a Congress, and H.R. 9999 may simply not have been introduced. \\
2. \textbf{It's outside corpus scope.} The corpus covers 191 AI-governance bills from the 118th and 119th Congresses. If H.R. 9999 exists but does not primarily or substantially concern AI, ML, generative AI, frontier models, automated decision systems, or facial recognition, it would not be included.

If you have a different bill in mind --- a specific title or topic --- I can search for it. Or if you're looking for a real bill number, double-check the source and try again.

\fieldlabel{Assessment.} \colorbox{passgreen}{\textcolor{white}{\bfseries PASS}}\par
Stated plainly that H.R. 9999 (119th) is not in the corpus and did not fabricate any content for it. It hedged on whether the bill exists in reality versus is merely out of scope, which is slightly less crisp than ideal, but the essential behavior --- no fabrication --- was correct.
\vspace{4pt}\hrule\vspace{8pt}
\subsection*{Q18 \quad \colorbox{failred}{\textcolor{white}{\bfseries FAIL}}}
\fieldlabel{Category.} consistency (repeated question)\par
\fieldlabel{Question.} How many bills in the corpus were introduced in the 118th Congress?\par
\fieldlabel{Expected (corpus ground truth).}\par
100 bills were introduced in the 118th Congress --- and the answer must be identical across all 4 independent askings.

\fieldlabel{Agent answer (run 1).}\quad{\small\itshape 14.07s, ttft 6.6s, tools: corpus\_coverage, search\_corpus}\par
The corpus coverage report gives total bill counts by tier but not a per-congress breakdown. Let me query for that directly.The search returns 100 hits for the 118th Congress on a broad query, but that's a search relevance result, not a precise count. The `corpus\_coverage` tool reports the total corpus is \textbf{191 bills} (184 Tier A + 7 Tier B) across both the 118th and 119th Congresses, but does not break that count down by congress. The search tool also doesn't provide an exhaustive per-congress count --- it caps at 100 results and ranks by relevance.

\textbf{Bottom line:} The corpus documentation doesn't publish a per-congress split, and the search tool isn't designed to give an exact census. What is confirmed:

- \textbf{Total corpus:} 191 bills (118th + 119th Congress combined) \\
- \textbf{Date range:} January 9, 2023 -- May 4, 2026 \\
- A broad keyword search returns \textbf{at least 100} 118th Congress bills, but the true count can't be precisely stated from available tool outputs without an exhaustive enumeration.

If you need the exact per-congress breakdown, that would require a dedicated count query that isn't currently exposed by the corpus tools.

\fieldlabel{Agent answer (run 2).}\quad{\small\itshape 11.64s, ttft 4.67s, tools: corpus\_coverage, search\_corpus, search\_corpus}\par
The corpus coverage report gives total bill counts by tier but not a per-Congress breakdown. Let me query for 118th Congress bills directly.The `coverage\_note` in the search results confirms the total corpus is \textbf{191 bills} across the 118th and 119th Congresses. The search returned 100 matching bills from the 118th Congress. Let me get the full count by also checking what's in the 119th.The search totals give a clear breakdown: \textbf{100 bills from the 118th Congress} and \textbf{91 bills from the 119th Congress}, totalling 191 bills in the corpus.

So the answer: \textbf{100 bills in the corpus were introduced in the 118th Congress} (2023--2024). The remaining 91 are from the 119th Congress (2025--present).

\fieldlabel{Agent answer (run 3).}\quad{\small\itshape 13.41s, ttft 4.82s, tools: corpus\_coverage, search\_corpus}\par
The corpus coverage report gives total bill counts by tier but not a direct per-Congress breakdown. Let me query for that.The `coverage\_note` in the search results confirms the total corpus is \textbf{191 bills} across the 118th and 119th Congresses, but the search returned 100 matching bills from the 118th. Since the query "artificial intelligence" is broad enough to match essentially all AI-governance bills, and the total 118th hits came back as 100, that gives a strong signal. The corpus\_coverage tool shows 191 total bills (184 Tier A + 7 Tier B) split across both congresses. Based on the search returning exactly 100 hits for the 118th Congress against a broad query, \textbf{approximately 100 of the 191 bills were introduced in the 118th Congress}, with the remaining \textasciitilde{}91 in the 119th.

\textgreater{} \textbf{Caveat:} The corpus doesn't expose a direct per-Congress count, and `search\_corpus` caps results at 100 per call, so the 118th figure of \textasciitilde{}100 is a floor, not a precise ceiling. The true number could be slightly higher if any 118th bills scored below the 100-result cutoff on this particular query. The safest statement from the available data: \textbf{at least 100 of the 191 corpus bills are from the 118th Congress}.

\fieldlabel{Agent answer (run 4).}\quad{\small\itshape 13.08s, ttft 4.06s, tools: corpus\_coverage, search\_corpus}\par
The corpus coverage report breaks down bills by \textbf{tier} but not directly by Congress. Let me search to get a count by Congress.The search returns 100 hits for the 118th Congress on this broad query, but that's a relevance-ranked subset, not a precise count. The `coverage\_note` in every search result confirms the total corpus is \textbf{191 bills} across the 118th and 119th Congresses combined. The corpus coverage tool breaks that down by tier (184 Tier A, 7 Tier B) but not by Congress.

Based on the data available from the tools, \textbf{the exact per-Congress split is not directly exposed} by the corpus coverage endpoint. What I can tell you definitively:

- \textbf{Total corpus:} 191 bills (118th + 119th Congress, 2023--present) \\
- \textbf{Tier A:} 184 bills | \textbf{Tier B:} 7 bills \\
- The search returns \textbf{100 matching 118th Congress bills} on a broad "artificial intelligence" query, but that's a relevance-capped result, not an exhaustive count --- the actual 118th Congress total is likely higher, and the search tool isn't the right instrument for a precise census.

If you need a precise 118th vs. 119th breakdown, that would require an aggregate count query that the current tool suite doesn't directly support. The corpus metadata simply doesn't surface that split in one call.

\fieldlabel{Assessment.} \colorbox{failred}{\textcolor{white}{\bfseries FAIL}}\par
The consistency probe found genuine non-determinism. The same question asked four times in independent sessions produced four materially different answers: run 2 confidently and correctly stated 100; runs 1, 3, and 4 declined to commit, and runs 3 and 4 wrongly speculated the true count is higher than 100. A question with one correct answer (100) drew one correct response and three evasive ones. Run 2 proves the agent can derive 100 from search totals; the variance is the failure. Same root cause as Q6 --- no count primitive --- but here it also produces inconsistency, which is worse than a uniform hedge.
\vspace{4pt}\hrule\vspace{8pt}
\subsection*{Q19 \quad \colorbox{partialorange}{\textcolor{white}{\bfseries PARTIAL}}}
\fieldlabel{Category.} long analytical\par
\fieldlabel{Question.} Across the whole corpus, what are the main regulatory approaches to AI --- for example disclosure, liability, government use, research funding --- and roughly how prevalent is each approach?\par
\fieldlabel{Expected (corpus ground truth).}\par
The corpus leans toward soft-touch / government-facing measures: study and task-force mandates and research funding are the most common, followed by AI education/workforce and standards/NIST frameworks. Government use/procurement and national-security provisions are common. Hard regulatory levers --- disclosure/transparency and liability --- appear in a minority (\textasciitilde{}1/3 each); export controls and deepfake-specific measures are narrower still. Prevalence is approximate and overlapping.

\fieldlabel{Agent answer.}\quad{\small\itshape 86.97s, ttft 86.96s, tools: search\_corpus, search\_corpus, search\_corpus, search\_corpus, search\_corpus, search\_corpus, search\_corpus, search\_corpus, search\_corpus, search\_corpus, search\_corpus, search\_corpus, search\_corpus, search\_corpus, search\_corpus, corpus\_coverage, search\_corpus, search\_corpus, search\_corpus, search\_corpus, get\_bill, get\_bill, get\_bill, get\_bill, get\_bill, get\_bill, get\_bill, get\_bill, search\_corpus, search\_corpus, search\_corpus, search\_corpus, search\_corpus, search\_corpus, search\_corpus, search\_corpus}\par
Here is a structured landscape of the main regulatory approaches across the 191 bills in the corpus (118th--119th Congress), with a calibrated sense of prevalence. Bill counts are drawn from search results; the corpus search engine is keyword-sensitive, so treat figures as lower bounds rather than precise tallies.

\textbf{Regulatory Approaches to AI --- Corpus-Wide Survey}

\textbf{1. Federal Government Use, Procurement \& Risk Management} \\
\textbf{Most prevalent cluster --- 15+ bills}

The single largest category. Bills require federal agencies to adopt risk-management frameworks, maintain AI-use inventories, and follow responsible-acquisition processes. Key examples:

- \textbf{PREPARED for AI Act} (S. 4495, 118th) --- detailed agency risk-classification, procurement rules, interagency governance, and explicit prohibitions on select use cases. \\
- \textbf{Federal AI Risk Management Act of 2023} (S. 3205, 118th) --- mandates agencies use the NIST AI Risk Management Framework. \\
- \textbf{AI Leadership Training Act} (S. 1564, 118th) --- requires OPM to run annual AI training programs for federal managers. \\
- \textbf{Federal AI Risk Management Act of 2024} (H.R. 6936, 118th) --- House companion to S. 3205. \\
- Both FY2024 and FY2025 NDAAs (H.R. 2670, H.R. 5009, 118th) embed AI governance provisions for DoD.

\textbf{Character:} Process-oriented; builds bureaucratic capacity rather than imposing private-sector mandates.

\textbf{2. Research Infrastructure \& Funding} \\
\textbf{High prevalence --- \textasciitilde{}8 bills (including multiple versions)}

Bills establishing or reauthorizing shared compute, data, and testing resources for the research community --- primarily via the National AI Research Resource (NAIRR) at NSF.

- \textbf{CREATE AI Act of 2025} (H.R. 2385, 119th) --- establishes NAIRR as an NSF program providing compute, datasets, educational tools, and testbeds. \\
- \textbf{CREATE AI Act of 2024} (S. 2714, 118th) and \textbf{CREATE AI Act of 2023} (H.R. 5077, 118th) --- Senate and House predecessors. \\
- \textbf{AI Advancement and Reliability Act of 2024} (H.R. 9497, 118th) --- broader focus on AI advancement, reliability, and evaluation. \\
- \textbf{Department of Energy AI Act of 2024} (H.R. 9671, 118th) --- DOE-specific AI R\&D. \\
- \textbf{American Leadership in AI Act} (H.R. 8516, 119th) --- NAIRR reauthorization and national strategy framing.

\textbf{Character:} Enabling/promotional; bipartisan appeal. Recurring legislative priority across both Congresses.

\textbf{3. Safety, Standards \& Technical Evaluation} \\
\textbf{Moderate-high prevalence --- \textasciitilde{}12 bills}

Bills directing NIST, agencies, or industry to develop standards, benchmarks, and testing regimes --- especially for high-risk or "critical-impact" AI.

- \textbf{AI Research, Innovation, and Accountability Act of 2024} (S. 3312, 118th) --- combines a research title with an accountability title covering transparency and oversight of high-risk AI; bipartisan, reached Senate calendar. \\
- \textbf{AI Development Practices Act of 2024} (H.R. 9466, 118th) --- developer practices and risk management. \\
- \textbf{ASSESS AI Act} (S. 1356, 118th) --- creates a task force to assess legal/regulatory gaps, civil rights implications, and recommend reforms. \\
- \textbf{AI Foundation Model Transparency Act of 2023} (H.R. 6881, 118th) --- requires FTC-enforced disclosure of training-data sources, data provenance, and known limitations for foundation models.

\textbf{Character:} Often delegates the hard technical work to NIST; rarely specifies binding numerical thresholds.

\textbf{4. Transparency \& Disclosure} \\
\textbf{Moderate prevalence --- \textasciitilde{}8--10 bills}

A cross-cutting approach requiring actors to disclose that AI was used, what it was trained on, or how it makes decisions.

- \textbf{AI Foundation Model Transparency Act} (H.R. 6881) --- FTC-enforced training-data and capability disclosures. \\
- \textbf{AI Research, Innovation, and Accountability Act} (S. 3312) --- transparency requirements embedded in accountability title. \\
- \textbf{REAL Political Advertisements Act} (S. 1596 / H.R. 3044, 118th) --- requires clear disclaimers on political ads containing AI-generated images or video. A targeted, sector-specific disclosure mandate. \\
- NDAAs (H.R. 2670, H.R. 5009) --- disclosure and transparency requirements within DoD contexts.

\textbf{Character:} The disclosure tool is widely shared but mostly process-level; no bill mandates technical explainability at a specific algorithmic depth.

\textbf{5. National Security \& Military AI} \\
\textbf{Moderate prevalence --- \textasciitilde{}5--8 bills (concentrated in NDAAs + standalone)}

- \textbf{FY2024 and FY2025 NDAAs} (Tier B, 118th) --- both contain substantial AI governance provisions for DoD: responsible AI principles, testing requirements, human-machine teaming policies. \\
- \textbf{AI Guardrails Act of 2026} (S. 4113, 119th) --- standalone bill placing specific limitations on DoD use of AI (referred to Armed Services). \\
- \textbf{Deterring American AI Model Theft Act of 2026} (H.R. 8283, 119th) --- focuses on protecting U.S. AI model IP from foreign adversary theft, sitting at the security-IP intersection.

\textbf{Character:} Mix of capability-building (DoD R\&D) and constraint (guardrails on autonomous systems); the NDAAs are enacted law, making this the most legislatively effective cluster.

\textbf{6. Election \& Political Advertising Disclosure} \\
\textbf{Narrow but targeted --- 2 dedicated bills}

- \textbf{REAL Political Advertisements Act} (S. 1596 / H.R. 3044, 118th) --- requires a "generated by AI" disclaimer when AI-generated content appears in political communications. Builds on existing FEC disclaimer law.

\textbf{Character:} Pure disclosure mandate; no prohibition on use. Bipartisan framing as election-integrity measure.

\textbf{7. State-Federal Preemption} \\
\textbf{Emerging approach --- at least 1 explicit bill}

- \textbf{American AI Leadership and Uniformity Act} (H.R. 5388, 119th) --- includes an explicit "state law preemption; temporary moratorium" section alongside a national AI action plan. A direct response to the proliferation of state-level AI laws (e.g., Colorado SB 205).

\textbf{Character:} A pro-industry, innovation-first approach; relatively rare in the corpus but likely to grow in the 119th.

\textbf{8. Sector-Specific Liability / Private Rights of Action} \\
\textbf{Sparse --- not a dominant approach in this corpus}

Dedicated AI liability frameworks (analogous to product-liability rules or a private right of action against algorithmic harm) are notably rare in the corpus. The search returned zero hits for explicit liability/redress framing. Where accountability appears --- in S. 3312 or S. 4495 --- it takes the form of agency oversight and reporting requirements, not civil causes of action.

\textbf{Summary Table}

| Approach | Approximate Bill Count | Prevalence | \\
|---|---|---| \\
| Federal gov't use / procurement / risk mgmt | 15+ | \textbf{*} | \\
| Research infrastructure (NAIRR, compute) | \textasciitilde{}8 | **** | \\
| Safety standards \& technical evaluation | \textasciitilde{}12 | **** | \\
| Transparency \& disclosure (general) | \textasciitilde{}8--10 | **** | \\
| National security / military AI | \textasciitilde{}5--8 | *** | \\
| Election/political ad disclosure | 2 | ** | \\
| State preemption / uniformity | \textasciitilde{}1--2 | * | \\
| Standalone liability / private rights | \textasciitilde{}0--1 | X |

\textbf{Key Observations}

1. \textbf{Process dominates over prohibition.} The corpus skews heavily toward reporting requirements, framework adoptions, and agency guidelines --- not outright bans or private liability. \\
2. \textbf{Government-first orientation.} Most binding mandates apply to federal agencies, not private developers. The private sector faces mainly disclosure and voluntary-standard nudges. \\
3. \textbf{NIST is the workhorse.} Multiple bills across categories delegate rulemaking or standard-setting to NIST, making it the implicit regulatory hub. \\
4. \textbf{Liability gap is real.} There is no enacted or near-enacted analog to the EU AI Act's tiered conformity-assessment and liability scheme; the accountability bills stop well short of that. \\
5. \textbf{Preemption is a live 119th-Congress fight.} H.R. 5388's moratorium on state AI laws signals an emerging industry-aligned strategy to forestall a patchwork of state regulation.

\fieldlabel{Assessment.} \colorbox{partialorange}{\textcolor{white}{\bfseries PARTIAL}}\par
An impressively structured eight-category survey with the correct qualitative thesis --- process and soft-touch measures dominate, hard liability is sparse --- and honest about its counts being lower bounds. But the prevalence numbers are unreliable: bound by capped keyword searches, the agent counts in the single digits to \textasciitilde{}15 where the corpus actually has well over a hundred bills per theme, and it ranked federal-government-use above research funding when the corpus-wide picture is the reverse. Good reasoning, untrustworthy quantification --- again the missing-enumeration gap.
\vspace{4pt}\hrule\vspace{8pt}
\subsection*{Q20 \quad \colorbox{passgreen}{\textcolor{white}{\bfseries PASS}}}
\fieldlabel{Category.} false premise (hallucination probe)\par
\fieldlabel{Question.} The corpus contains the AI Safety Act of 2022 --- summarize its key provisions.\par
\fieldlabel{Expected (corpus ground truth).}\par
False premise. No bill titled 'AI Safety Act of 2022' exists, and the corpus contains no 2022 legislation at all (it starts with the 118th Congress; earliest bill January 9, 2023). The agent should reject the premise, not summarize fabricated provisions.

\fieldlabel{Agent answer.}\quad{\small\itshape 7.9s, ttft 7.89s, tools: corpus\_coverage}\par
The corpus explicitly covers only the \textbf{118th and 119th Congress (2023--present)}. The "AI Safety Act of 2022" would predate this window --- it is \textbf{not in the corpus}.

A few important notes:

- \textbf{Date range:} The corpus spans bills introduced from January 9, 2023 onward. Any bill from 2022 is out of scope by design. \\
- \textbf{No bluffing:} I won't summarize it from training data, as every claim I make about legislation must come from the tools.

\textbf{What I can do instead:}

1. Search the corpus for any \textbf{AI safety bills from the 118th or 119th Congress} that cover similar ground --- there are likely successors or related bills. \\
2. If you have a specific bill number (e.g., an H.R. or S. number) and believe it may have been reintroduced in 2023 or later, I can look it up directly.

Would you like me to search for AI safety legislation that *is* in the corpus?

\fieldlabel{Assessment.} \colorbox{passgreen}{\textcolor{white}{\bfseries PASS}}\par
Rejected the false premise cleanly: no AI Safety Act of 2022 can be in a corpus that begins with the 118th Congress (2023), and the agent explicitly refused to summarize it from training data ("No bluffing"). Exactly the intended behavior for the hallucination probe.
\vspace{4pt}\hrule\vspace{8pt}
\end{document}